\chapter{System Design}\label{ch:iv:system-design}

``Design a matching engine.'' Forty-five minutes, a whiteboard, no other words. The candidate who spends five minutes asking
how many instruments, how many messages a second at peak, what latency at the ninety-ninth percentile, and what happens when
the process dies has already done better than the one who began with the order book's data structure. \omterm{def:iv:system-design:interview}{System-design interviews}
at trading firms ask for the systems the series builds: a matching engine and exchange simulator (One Quant Book~10, chapter~26),
feed handlers and order gateways (Book~13), market-data distribution, backtesters (Book~7 and Book~15) and risk services. What is
scored is the method and the numbers, not a single right architecture.

\section{The method}

\begin{definition}[System-design interview, capacity estimate]\label{def:iv:system-design:interview}
A \emph{system-design interview}\index{system-design interview} asks the candidate to design a system from a one-line
description within a fixed time, and assesses how requirements are established, how the design is sized, how its components
and interfaces are chosen, and how it fails. A \emph{capacity estimate}\index{capacity estimate} is the back-of-envelope
calculation of the load a design must carry (messages a second at peak, bytes a day, memory, connections) from stated
assumptions, used to choose between designs.
\end{definition}

\begin{method}[Answering a design question]\label{met:iv:system-design:method}
\begin{enumerate}
\item \emph{Requirements}: functions (what it must do), scale (instruments, users, messages a second at peak and on average),
latency (median and tail), correctness (ordering, exactly-once), availability (what may be lost on failure, how fast to recover).
\item \emph{Capacity}: a \omterm{def:iv:system-design:interview}{capacity estimate} for messages, bandwidth, memory and storage, with the peak-to-average ratio stated.
\item \emph{Interfaces}: the messages in and out, and their sequencing.
\item \emph{Core design}: the data structures and the threading model of the critical path first.
\item \emph{Failure}: what happens when each component dies, when a consumer is slow, when the network drops packets.
\item \emph{Numbers again}: check the design against the \omterm{def:iv:system-design:interview}{capacity estimate} and the latency budget.
\end{enumerate}
\end{method}

\begin{example}[Sizing a feed]\label{ex:iv:system-design:feed}
Suppose (assumptions to be stated as such) 200\,000 messages a second on average over a 6.5-hour session and 64 bytes a
message: $200\,000 \times 23\,400 \times 64 \approx 300$ gigabytes a day to capture. At a peak of one million messages a second
the payload alone is 512 megabits a second, and with 42 bytes of Ethernet, IPv4 and UDP headers ($14 + 20 + 8$) per message sent singly,
848 megabits: comfortably inside 10 gigabits, but not inside 1.
\end{example}

\section{The order book and the matching engine}

The matching engine is the question asked most. Its core is a price-time priority book: price levels on each side in sorted
order, a first-in-first-out queue of orders within each level, and an index from order identifier to its position so that a
cancel costs constant time (\cref{lst:iv:system-design:match}). Around the core sit the pieces that make it an exchange: a
sequencer that stamps every inbound message with a single order, a journal that writes the sequenced input to durable storage
before or while it is processed, the market-data publisher and the order-entry gateways (\cref{fig:iv:system-design:engine}).

\omcode{code/interviews/26-system-design/cpp/iv_book.hpp}{37}{56}{The matching loop of a price-time book: consume the best
opposite level in time order while the incoming order crosses. The chapter's test compares every fill and best price with a
brute-force book on 200 random order streams.}\label{lst:iv:system-design:match}

\begin{omfigure}
\begin{tikzpicture}[font=\footnotesize,>={Stealth},
  box/.style={draw=omDef,rounded corners=2pt,fill=omDef!6,minimum height=0.8cm,text width=2.2cm,align=center,inner sep=2pt},
  alt/.style={draw=omThm,rounded corners=2pt,fill=omThm!8,minimum height=0.8cm,text width=2.2cm,align=center,inner sep=2pt}]
\node[box] (gw) at (0,0) {order-entry gateways};
\node[box] (seq) at (3.1,0) {sequencer};
\node[box] (me) at (6.2,0) {matching engine (one thread per partition)};
\node[box] (md) at (9.3,0.9) {market-data publisher};
\node[box] (acks) at (9.3,-0.9) {acknow\-ledge\-ments and fills};
\node[alt] (jr) at (3.1,-1.7) {journal of the sequenced input};
\node[alt] (rep) at (6.2,-1.7) {replica replaying the journal};
\draw[->,thick,omDef] (gw) -- (seq);
\draw[->,thick,omDef] (seq) -- (me);
\draw[->,thick,omDef] (me) -- (md);
\draw[->,thick,omDef] (me) -- (acks);
\draw[->,thick,omThm] (seq) -- (jr);
\draw[->,thick,omThm] (jr) -- (rep);
\draw[->,thick,omDef] (acks.south) -- ++(0,-1.75) -| (gw.south);
\node[font=\scriptsize,anchor=north] at (4.6,-3.2) {acknowledgements go back to the participants through the gateways};
\end{tikzpicture}

\omcaption{A matching engine as a replicated state machine: the sequencer gives every inbound message one order and journals
it; the engine is a deterministic function of the sequenced input, so a replica replaying the journal reaches the same state
and can take over (One Quant Book~13, chapter~24).}\label{fig:iv:system-design:engine}
\end{omfigure}

Determinism is the design's key property: if the engine's output depends only on its sequenced input (no clocks read, no hash-map
iteration order, no thread interleavings inside a partition), recovery is replay and a standby replica is a copy. Scaling is by
partition: instruments are divided among engine threads or machines, each with its own book, which works because orders on
different instruments do not interact (a spread order across two partitions needs a coordinating component).

\section{Market data and the backtester}

A market-data service fans out one feed to many consumers. Its design questions are sequencing (every message numbered per
channel), recovery (a snapshot channel and a retransmission service, One Quant Book~10, chapter~26), and slow consumers, which must
never slow the publisher: each consumer gets its own queue, and a consumer that falls too far behind is disconnected or switched to
snapshots (\cref{fig:iv:system-design:md}).

\begin{omfigure}
\begin{tikzpicture}[font=\footnotesize,>={Stealth},
  box/.style={draw=omDef,rounded corners=2pt,fill=omDef!6,minimum height=0.8cm,text width=2.3cm,align=center,inner sep=2pt},
  alt/.style={draw=omThm,rounded corners=2pt,fill=omThm!8,minimum height=0.8cm,text width=2.3cm,align=center,inner sep=2pt}]
\node[box] (fh) at (0,0) {feed handler (arbitrates lines A and B)};
\node[box] (bk) at (3.3,0) {book builder};
\node[box] (pub) at (6.6,0) {publisher: incremental channel};
\node[alt] (snap) at (6.6,-1.7) {snapshot and retransmission service};
\node[box] (c1) at (10,1.4) {consumer: strategy};
\node[box] (c2) at (10,0.2) {consumer: risk};
\node[alt] (c3) at (10,-1.35) {slow consumer: own queue, falls back to snapshots};
\draw[->,thick,omDef] (fh) -- (bk);
\draw[->,thick,omDef] (bk) -- (pub);
\draw[->,thick,omDef] (pub) -- (c1.west);
\draw[->,thick,omDef] (pub) -- (c2.west);
\draw[->,thick,omDef] (pub) -- (c3.west);
\draw[->,thick,omThm] (bk) |- (snap);
\draw[->,thick,omThm,dashed] (snap.east) -- (c3.west |- snap.east);
\end{tikzpicture}

\omcaption{A market-data service: sequenced incremental updates for everyone, a snapshot and retransmission path for
consumers that miss messages or fall behind, and a separate queue per consumer so that none can slow the
publisher.}\label{fig:iv:system-design:md}
\end{omfigure}

A backtester (One Quant Book~7, chapter~17, and One Quant Book~15, chapter~11) is an event loop over a merged stream of recorded
events, with a simulated clock, a fill model for orders against the recorded book, latencies applied to both directions, and
the strategy code called through the same interface it uses in production. Its design questions are reproducibility (seeded
randomness, versioned data and code), realism (queue position, latency, fees) and speed (partitioning by day or instrument).

\section{The risk service and the question of state}

A pre-trade risk service sits on the order path (One Quant Book~11, chapter~27): every order is checked against limits (size,
price collar, position, credit) before it leaves, so its state (positions, open orders, limits) must be current and its latency
small. The interview questions are where the state lives (in the gateway's process for speed, in a central service for a firm-wide
view, usually both with the central view as the authority), how it is kept consistent (the gateway reserves against a limit
before sending and releases on cancel or fill), and what happens when it fails: a risk check that cannot run must block orders
(fail closed), and a kill switch must work when everything else does not.

\section{Worked answers}

\begin{example}[A latency budget]\label{ex:iv:system-design:budget}
\emph{``The desk wants tick-to-order under 10 microseconds at the 99th percentile. Split the budget.''} Write the path and
give each stage a number: network card to application 1.5 (kernel bypass), feed decode 0.5, book update 1.0, strategy 3.0,
pre-trade risk 1.0, order encode 0.5, application to wire 2.0: 9.5 microseconds, leaving 0.5 of margin. Then the point the
interviewer is waiting for: these are percentiles, and percentiles do not add. If each of seven stages independently has a
1\% chance of a slow path, the chance that at least one is slow on a given message is $1 - 0.99^7 \approx 6.8\%$, so a
budget of per-stage 99th percentiles says little about the end-to-end 99th percentile. Measure end to end, on the real
path, with hardware timestamps, and budget the tail (cache misses, interrupts, a page fault) explicitly
(\cref{ch:iv:concurrency-operating-systems-and-networks}).
\end{example}

\section{Question bank}

\begin{interviewq}[$\star$ \iqroles{developer} \iqfirm{market maker}]\label{iq:iv:system-design:1}
A feed averages 200\,000 messages a second over a 6.5-hour session, 64 bytes each. How much storage does a day of capture take?
\end{interviewq}

\begin{interviewq}[$\star$ \iqroles{developer} \iqfirm{proprietary firm}]\label{iq:iv:system-design:2}
At a peak of one million 64-byte messages a second, what bandwidth does the feed need? Is a 1-gigabit link enough? A 10-gigabit
one?
\end{interviewq}

\begin{interviewq}[$\star$ \iqroles{developer} \iqfirm{crypto firm}]\label{iq:iv:system-design:3}
You are asked to design a matching engine for a new venue. List the questions you ask in the first five minutes.
\end{interviewq}

\begin{interviewq}[$\star$ \iqroles{developer} \iqfirm{market maker}]\label{iq:iv:system-design:4}
Which operations must an order book support, and with which complexities, for a matching engine?
\end{interviewq}

\begin{interviewq}[$\star\star$ \iqroles{developer} \iqfirm{market maker}]\label{iq:iv:system-design:5}
Estimate the memory for full-depth books of 10\,000 instruments with five million live orders, 48 bytes an order, and up to 1\,000
price levels a side at 16 bytes a level.
\end{interviewq}

\begin{interviewq}[$\star\star$ \iqroles{developer} \iqfirm{crypto firm}]\label{iq:iv:system-design:6}
How do you make a matching engine recoverable after a crash with no lost or duplicated orders? What must be true of the engine's
code for your answer to work?
\end{interviewq}

\begin{interviewq}[$\star\star$ \iqroles{developer} \iqfirm{proprietary firm}]\label{iq:iv:system-design:7}
Design a service that distributes one market-data feed to fifty internal consumers, some of them slow.
\end{interviewq}

\begin{interviewq}[$\star\star$ \iqroles{developer, researcher} \iqfirm{systematic fund}]\label{iq:iv:system-design:8}
Design a backtester for intraday strategies on recorded order-book data. What makes its results reproducible and believable?
\end{interviewq}

\begin{interviewq}[$\star\star$ \iqroles{developer} \iqfirm{market maker}]\label{iq:iv:system-design:9}
Write the matching loop of a price-time priority book for limit orders, and say how you would test it.
\end{interviewq}

\begin{interviewq}[$\star\star\star$ \iqroles{developer, risk} \iqfirm{bank}]\label{iq:iv:system-design:10}
Design the pre-trade risk check for a firm with order gateways in three data centres. Where does the state live, how is it kept
consistent, and what happens when the risk service is down?
\end{interviewq}

\begin{interviewq}[$\star\star\star$ \iqroles{developer} \iqfirm{crypto firm}]\label{iq:iv:system-design:11}
Your matching engine handles 10\,000 instruments; the peak is three million messages a second and one thread handles one million.
How many partitions do you run at no more than half load, and what breaks when an order spans two partitions?
\end{interviewq}

\begin{interviewq}[$\star\star\star$ \iqroles{developer} \iqfirm{proprietary firm}]\label{iq:iv:system-design:12}
Design failover for the matching engine so that a standby takes over within a second of the primary's death with no lost
acknowledged order.
\end{interviewq}

\begin{interviewq}[$\star\star\star$ \iqroles{developer, risk} \iqfirm{multi-manager fund}]\label{iq:iv:system-design:13}
A real-time risk service revalues 100\,000 accounts of 50 positions each; on each of 1\,000 price ticks a second, about 1\% of all
positions are affected. How many position updates a second is that, and how do you design for it?
\end{interviewq}

\begin{omsources}
\item One Quant Book~7, chapter~17; One Quant Book~10, chapters 1 and~26; One Quant Book~11, chapter~27; One Quant Book~13, chapters
18--24; One Quant Book~15, chapter~11.
\item The series' exchange simulator, \texttt{firm.exchsim}, and its protocol document (\texttt{code/firm/exchsim/PROTOCOL.md}), as a
worked design of a venue.
\end{omsources}
